\cotefile{136}
\scanpage{220}
% slanted
\margin{De Rham P.D.}

\textcircled{3} Sur $\Gamma^*M\otimes\Lambda^*M$ \deleted{(algèbre bigraduée)} il y a un seul op. $d$ qui soit
\deleted{a) $A$ linéaire} a) une \inserted{anti}dérivation si \deleted{$\Gamma^*(M)$ de degrés pairs et $\Lambda^1M$ de degré impair}
\inserted{$A$-algèbre graduée,} b) \inserted{antidérivation de degré 1} compatible avec les puiss. divisées sur $(\Gamma^*M)^+$ i.e. $d(x^{(i)})=x^{(i-1)}\,dx$, et
\unsure{c)} $d$ est nulle sur $\Lambda^1M$, et sur $\Gamma^1 M$ c'est l'isom can $\Gamma^1 M\longrightarrow\Lambda^1 M$.
Alors $d$ satisfait $d^2=0$, et applique $\Gamma^p\otimes\Lambda^q$ dans $\Gamma^{p\unsure{-|+}1}\otimes\Lambda^{q+1}$. Il est $\Lambda^*M$-linéaire.
C'est aussi \unsure{du|dit|le} ``complexe de De Rham à puissances divisées'' de la $A$-algèbre à p.d. $\Gamma^*M$ (avec son idéal $\Gamma^+M$).
\margin{? $\big\{$}
Si $M$ est plat sur $A$, alors on a (\unsure{graduant|gradué} par le degré de $\Lambda^*$)
\[
 H^{\unsure{*|i}}_{\mathrm{DR\,p.d.}}(\Gamma^*(M)\otimes\Lambda^* M)
 \simeq\begin{cases}0&\text{si }i\ne0\\
 \underset{\substack{\rotatebox{90}{$\simeq$}\\\Gamma^0\otimes\Lambda^0}}{A}&\text{si }i=0\end{cases}
\]

% slanted
\margin{Koszul P.D.}

\textcircled{4} Sur $\Gamma^*M\otimes\Lambda^*M$, il y a un seul op. $\delta$ qui soit
a) dérivation de \inserted{$A$-}algèbre graduée b) compatible avec p.d. sur $\Gamma^+(M)$ ($d(x^{(i)})=x^{(i-1)}dx$ si $x\in\Gamma^+$).
c) nulle sur $\Gamma^1 M$ \deleted{(donc sur $\Gamma^*M$ \unsure{entier})} et induisant sur $\Lambda^1M$ l'isom $\Lambda^1M\longrightarrow\Gamma^1M$.
\deleted{Si $M$ \unsure{libre} \inserted{plat} sur $A$, et si $A$ de \unsure{caractéristique}} \inserted{Alors $\delta^2=0$, il applique $\Gamma^p\otimes\Lambda^q$ dans $\Gamma^{p+1}\otimes\Lambda^{q-1}$,} \inserted{est $\Gamma^*M$-linéaire.}
Ecrivant les degrés en bas, si $M$ est plat \inserted{(inutile)} sur $A$ \deleted{\unsure{$M$}} de car. $0$, on a
\[
 H_i^{\mathrm{Kos\,p.d.}}(\Gamma^*M\otimes\Lambda^*M)\simeq
 \begin{cases}0&\text{si }i\ne0\\
 \underset{\substack{\rotatebox{90}{$\simeq$}\\\Gamma^0\otimes\Lambda^0}}{A}&\text{si }i=0\end{cases}
\]

Si $A$ de car $p>0$, $M$ plat, on doit avoir une histoire d'\unsure{opérations|opération} de Cartier\unsure{...}.

% archive sheet number: 219

